% GENERATED FILE. Edit canonical lessons, then run npm run generate:books.
% canonical-source-hash: fnv1a64:f180dea54d916f1b
% canonical-lessons: ZH-C172-weidao, ZH-C172-suan, ZH-C172-ku, ZH-C172-xiang, ZH-C172-dan

\chapter{Sweet and Sour}
\label{ch:zh-sweet-and-sour}
\hlchaptermodality{172}

\begin{chapteropening}
\textbf{By the end of this chapter:} \emph{I can say how food tastes: good, sour, bitter, fragrant or bland.}

Complete the last lesson of chapter 172: I can say how food tastes: good, sour, bitter, fragrant or bland.
\end{chapteropening}

\section[wèidao]{wèidao — taste; flavour}

\label{lesson:ZH-C172-weidao}



\begin{quote}
Before the new one: say the Mandarin for to press a button, then the Mandarin for to touch.
\end{quote}

\subsection*{You'll want to know: wèidao}
\emph{wèidao} — \textquotedblleft{}taste; flavour\textquotedblright{}.

\begin{quote}
\emph{zhège cài de wèidao hěn hǎo} — This dish tastes good.
\end{quote}

Its characters come later; for now it is a word for the ear.

\begin{cousinweb}
\emph{wèidao} is built from \emph{wèi} 'taste' and \emph{dao}, both inherited from Middle Chinese, \emph{mj+jH} and \emph{dawX} (Guangyun, in Baxter's transcription). The dictionary glosses the first as the five tastes, sour, salty, sweet, bitter and hot, and the second 'a road'. In the word the second loses its tone.
\end{cousinweb}

\begin{grammarlens}[title={What you've built}]
One more word for asking for what I want.
\end{grammarlens}

\subsection*{Your turn}
\begin{itemize}
\raggedright
  \item \emph{Say it:} \emph{wèidao}
  \item \emph{Say it:} \emph{wèidao}, once more
  \item \emph{Say it:} \emph{zhège cài de wèidao hěn hǎo} — This dish tastes good.
  \item \emph{Recall:} say the Mandarin for to press a button, then the Mandarin for to touch, then say \emph{wèidao} again
\end{itemize}

\begin{tcolorbox}[breakable,colback=teal!4,colframe=teal!35!black,title={Before you move on}]
What does wèidao mean? (\textquotedblleft{}Taste; flavour\textquotedblright{}.) Say it once more.
\end{tcolorbox}
\section[suān]{suān — sour}

\label{lesson:ZH-C172-suan}



\begin{quote}
Before the new one: say the Mandarin for to touch, then the Mandarin for a flavour.
\end{quote}

\subsection*{You'll want to know: suān}
\emph{suān} — \textquotedblleft{}sour\textquotedblright{}.

\begin{quote}
\emph{zhège píngguǒ hěn suān} — This apple is sour.
\end{quote}

Its characters come later; for now it is a word for the ear.

\begin{cousinweb}
\emph{suān} is inherited from Middle Chinese \emph{swan} (Guangyun, in Baxter's transcription), glossed 'vinegar': the taste of vinegar.
\end{cousinweb}

\begin{grammarlens}[title={What you've built}]
One more word for asking for what I want.
\end{grammarlens}

\subsection*{Your turn}
\begin{itemize}
\raggedright
  \item \emph{Say it:} \emph{suān}
  \item \emph{Say it:} \emph{suān}, once more
  \item \emph{Say it:} \emph{zhège píngguǒ hěn suān} — This apple is sour.
  \item \emph{Recall:} say the Mandarin for to touch, then the Mandarin for a flavour, then say \emph{suān} again
\end{itemize}

\begin{tcolorbox}[breakable,colback=teal!4,colframe=teal!35!black,title={Before you move on}]
What does suān mean? (\textquotedblleft{}Sour\textquotedblright{}.) Say it once more.
\end{tcolorbox}
\section[kǔ]{kǔ — bitter}

\label{lesson:ZH-C172-ku}



\begin{quote}
Before the new one: say the Mandarin for a flavour, then the Mandarin for sour.
\end{quote}

\subsection*{You'll want to know: kǔ}
\emph{kǔ} — \textquotedblleft{}bitter\textquotedblright{}.

\begin{quote}
\emph{zhège kāfēi tài kǔ le} — This coffee is too bitter.
\end{quote}

Its characters come later; for now it is a word for the ear.

\begin{cousinweb}
\emph{kǔ} is inherited from Middle Chinese \emph{khuX} (Guangyun, in Baxter's transcription), glossed 'coarse', 'toil' and 'suffering'; the taste itself is not given there in so many words.
\end{cousinweb}

\begin{grammarlens}[title={What you've built}]
One more word for asking for what I want.
\end{grammarlens}

\subsection*{Your turn}
\begin{itemize}
\raggedright
  \item \emph{Say it:} \emph{kǔ}
  \item \emph{Say it:} \emph{kǔ}, once more
  \item \emph{Say it:} \emph{zhège kāfēi tài kǔ le} — This coffee is too bitter.
  \item \emph{Recall:} say the Mandarin for a flavour, then the Mandarin for sour, then say \emph{kǔ} again
\end{itemize}

\begin{tcolorbox}[breakable,colback=teal!4,colframe=teal!35!black,title={Before you move on}]
What does kǔ mean? (\textquotedblleft{}Bitter\textquotedblright{}.) Say it once more.
\end{tcolorbox}
\section[xiāng]{xiāng — smelling good; fragrant; tasty}

\label{lesson:ZH-C172-xiang}



\begin{quote}
Before the new one: say the Mandarin for sour, then the Mandarin for bitter.
\end{quote}

\subsection*{You'll want to know: xiāng}
\emph{xiāng} — \textquotedblleft{}smelling good; fragrant; tasty\textquotedblright{}.

\begin{quote}
\emph{nǐ zuò de fàn zhēn xiāng} — The food you cook smells wonderful.
\end{quote}

Its characters come later; for now it is a word for the ear.

\begin{cousinweb}
\emph{xiāng} is inherited from Middle Chinese \emph{xjang} (Guangyun, in Baxter's transcription), glossed 'fragrant'.
\end{cousinweb}

\begin{grammarlens}[title={What you've built}]
One more word for asking for what I want.
\end{grammarlens}

\subsection*{Your turn}
\begin{itemize}
\raggedright
  \item \emph{Say it:} \emph{xiāng}
  \item \emph{Say it:} \emph{xiāng}, once more
  \item \emph{Say it:} \emph{nǐ zuò de fàn zhēn xiāng} — The food you cook smells wonderful.
  \item \emph{Recall:} say the Mandarin for sour, then the Mandarin for bitter, then say \emph{xiāng} again
\end{itemize}

\begin{tcolorbox}[breakable,colback=teal!4,colframe=teal!35!black,title={Before you move on}]
What does xiāng mean? (\textquotedblleft{}Smelling good; fragrant; tasty\textquotedblright{}.) Say it once more.
\end{tcolorbox}
\section[dàn]{dàn — bland; weak (in taste)}

\label{lesson:ZH-C172-dan}



\begin{quote}
Before the new one: say the Mandarin for bitter, then the Mandarin for smelling good.
\end{quote}

\subsection*{You'll want to know: dàn}
\emph{dàn} — \textquotedblleft{}bland; weak (in taste)\textquotedblright{}.

\begin{quote}
\emph{tāng tài dàn le} — The soup is too bland.
\end{quote}

Its characters come later; for now it is a word for the ear.

\begin{cousinweb}
\emph{dàn} is inherited from Middle Chinese \emph{damX} (Guangyun, in Baxter's transcription), glossed 'a thin taste'.
\end{cousinweb}

\begin{grammarlens}[title={What you've built}]
One more word for asking for what I want.
\end{grammarlens}

\subsection*{Your turn}
\begin{itemize}
\raggedright
  \item \emph{Say it:} \emph{dàn}
  \item \emph{Say it:} \emph{dàn}, once more
  \item \emph{Say it:} \emph{tāng tài dàn le} — The soup is too bland.
  \item \emph{Recall:} say the Mandarin for bitter, then the Mandarin for smelling good, then say \emph{dàn} again
\end{itemize}

\begin{tcolorbox}[breakable,colback=teal!4,colframe=teal!35!black,title={Before you move on}]
What does dàn mean? (\textquotedblleft{}Bland; weak (in taste)\textquotedblright{}.) Say it once more.
\end{tcolorbox}
